\chapter{1997年上海卷（理）}
\section{单选题}

\begin{problem}
设全集是实数集 \(\R\)，\(M=\left\{x \mid x \leqslant 1+\sqrt{2}, x \in \R\right\}\)，\(N=\{1,2,3,4\}\). 则 \(\overline{M} \cap N\) 等于……（\quad）

\choices
    {\(\{4\}\).}
    {\(\{3,4\}\).}
    {\(\{2,3,4\}\).}
    {\(\{1,2,3,4\}\).}

\end{problem}

\begin{problem}
三个数 \(6^{0.7}\)，\(0.7^{6}\)，\(\log_{0.7}6\) 的大小顺序是……（\quad）

\choices
    {\(0.7^6 < \log_{0.7}6 < 6^{0.7}\).}
    {\(0.7^6 < 6^{0.7} < \log_{0.7}6\).}
    {\(\log_{0.7}6 < 6^{0.7} < 0.7^6\).}
    {\(\log_{0.7}6 < 0.7^6 < 6^{0.7}\).}

\end{problem}

\begin{problem}
在下列命题中，假命题是……（\quad）

\choices
    {若 \(a\)，\(b\) 是异面直线，则一定存在平面 \(\alpha\) 过 \(a\) 且与 \(b\) 平行.}
    {若 \(a\)，\(b\) 是异面直线，则一定存在平面 \(\alpha\) 过 \(a\) 且与 \(b\) 垂直.}
    {若 \(a\)，\(b\) 是异面直线，则一定存在平面 \(\alpha\) 与 \(a\)，\(b\) 所成角相等.}
    {若 \(a\)，\(b\) 是异面直线，则一定存在平面 \(\alpha\) 与 \(a\)，\(b\) 的距离相等.}

\end{problem}

\begin{problem}
设 \(\theta \in \left(\frac{3\pi}{4}, \pi\right)\)，则关于 \(x\)，\(y\) 的方程 \(x^2 \csc \theta - y^2 \sec \theta = 1\) 所表示的曲线是……（\quad）

\choices
    {实轴在 \(y\) 轴上的双曲线.}
    {实轴在 \(x\) 轴上的双曲线.}
    {长轴在 \(y\) 轴上的椭圆.}
    {长轴在 \(x\) 轴上的椭圆.}

\end{problem}

\begin{problem}
如果直线 \(l\) 沿 \(x\) 轴负方向平移 \(3\) 个单位，再沿 \(y\) 轴正方向平移 \(1\) 个单位后，又回到原来的位置，那么直线 \(l\) 的斜率是……（\quad）

\choices
    {\(-\frac{1}{3}\).}
    {\(-3\).}
    {\(\frac{1}{3}\).}
    {\(3\).}

\end{problem}

\begin{problem}
设 \(f(n)=\frac{1}{n+1}+\frac{1}{n+2}+\frac{1}{n+3}+\cdots+\frac{1}{2n}\)（\(n \in \mathbb{N}\)），那么 \(f(n+1)-f(n)\) 等于（\quad）

\choices
    {\(\frac{1}{2n+1}\).}
    {\(\frac{1}{2n+2}\).}
    {\(\frac{1}{2n+1}+\frac{1}{2n+2}\).}
    {\(\frac{1}{2n+1}-\frac{1}{2n+2}\).}

\end{problem}

\section{填空题}

\begin{problem}
二次曲线 \(\left\{ \begin{array}{l} x = 5\cos \theta \\ y = 3\sin \theta \end{array} \right.\)（\(\theta\) 是参数）的左焦点坐标是 \fillinblank{}.

\end{problem}

\begin{answer}
\((-4,0)\)
\end{answer}

\begin{problem}
方程 \(\sin 2x = \frac{1}{2}\) 在 \(\left[-2\pi, 2\pi\right]\) 内解的个数是 \fillinblank{}.

\end{problem}

\begin{answer}
8
\end{answer}

\begin{problem}
已知 \(a=\frac{-3-\i}{1+2\i}\)（\(\i\) 是虚数单位），那么 \(a^4=\fillinblank{}\).

\end{problem}

\begin{answer}
\(-4\)
\end{answer}

\begin{problem}
设圆 \(x^2+y^2-4x-5=0\) 的弦 \(AB\) 的中点为 \(P(3,1)\)，则直线 \(AB\) 的方程是 \fillinblank{}.

\end{problem}

\begin{answer}
\(x+y-4=0\)
\end{answer}

\begin{problem}
若 \((3x+1)^n\)（\(n \in \mathbb{N}\)）的展开式中各项系数的和是 \(256\)，则展开式中 \(x^2\) 的系数是 \fillinblank{}.

\end{problem}

\begin{answer}
10
\end{answer}

\begin{problem}
函数 \(f(x)=3\sin x\cos x-4\cos^2 x\) 的最大值是 \fillinblank{}.

\end{problem}

\begin{answer}
\(\frac{1}{2}\)
\end{answer}

\section{选作题}

\begin{problem}
（\(A_1\)）设 \(0<a<b\)，则 \(\displaystyle\lim_{n \to \infty}\frac{4b^n}{a^n-b^n}=\fillinblank{}\).

\end{problem}

\begin{answer}
\(-4\)
\end{answer}

\begin{problem}
（\(A_1\)）设正四棱锥底面边长为 \(4\mathrm{cm}\)，侧面和底面所成的二面角是 \(60^\circ\)，则这个棱锥的侧面积是 \(\fillinblank{}\mathrm{cm}^2\).

\end{problem}

\begin{answer}
32
\end{answer}

\begin{problem}
（\(A_1\)）圆柱形容器的内壁底半径为 \(5\mathrm{cm}\)，两个直径为 \(5\mathrm{cm}\) 的玻璃小球都浸没于容器的水中. 若取出这两个小球，则容器内的水面将下降 \fillinblank{} \(\mathrm{cm}\).

\end{problem}

\begin{answer}
\(\frac{5}{3}\)
\end{answer}

\begin{problem}
（\(A_1\)）从集合 \(\{0,1,2,3,5,7,11\}\) 中任取 \(3\) 个元素分别作为直线方程 \(Ax+By+C=0\) 中的 \(A\)，\(B\)，\(C\)，所得恰好经过坐标原点的直线有 \fillinblank{} 条（结果用数值表示）.

\end{problem}

\begin{answer}
30
\end{answer}

\begin{problem}
（\(B_1\)）\(\displaystyle\lim_{n \to \infty}\left(1-\frac{2}{n}\right)^{-2n}=\fillinblank{}\).

\end{problem}

\begin{answer}
\(\e^4\)
\end{answer}

\begin{problem}
（\(B_1\)）设 \(\overrightarrow{a}=(m+1)\overrightarrow{i}-3\overrightarrow{j}\)，\(\overrightarrow{b}=\overrightarrow{i}+(m-1)\overrightarrow{j}\)，\(\left(\overrightarrow{a}+\overrightarrow{b}\right)\perp\left(\overrightarrow{a}-\overrightarrow{b}\right)\)，则 \(m=\fillinblank{}\).

\end{problem}

\begin{answer}
\(-2\)
\end{answer}

\begin{problem}
（\(B_1\)）设 \(f(x)=x^2(2-x)\)，则 \(f(x)\) 的单调递增区间是 \fillinblank{}.

\end{problem}

\begin{answer}
\(\left(0,\frac{4}{3}\right)\)
\end{answer}

\begin{problem}
（\(B_1\)）从集合 \(\{0,1,2,3,5,7,11\}\) 中任取 \(3\) 个元素分别作为直线方程 \(Ax+By+C=0\) 中的 \(A\)，\(B\)，\(C\)，所得恰好经过坐标原点的直线的概率是 \fillinblank{}.

\end{problem}

\begin{answer}
\(\frac{1}{7}\)
\end{answer}

\section{解答题}

\begin{problem}
已知 \(\tan \frac{\alpha}{2}=\frac{1}{2}\)，求 \(\sin \left(\alpha+\frac{\pi}{6}\right)\) 的值.

\end{problem}

\begin{solution}
解法一：因为
\[
\sin \alpha=\frac{2 \tan \frac{\alpha}{2}}{1+\tan^2 \frac{\alpha}{2}}=\frac{2 \times \frac{1}{2}}{1+\frac{1}{4}}=\frac{4}{5}
\]
\[
\cos \alpha=\frac{1-\tan^2 \frac{\alpha}{2}}{1+\tan^2 \frac{\alpha}{2}}=\frac{1-\frac{1}{4}}{1+\frac{1}{4}}=\frac{3}{5}
\]
因此 \(\sin\left(\alpha+\frac{\pi}{6}\right)=\sin \alpha \cos \frac{\pi}{6}+\cos \alpha \sin \frac{\pi}{6}=\frac{4}{5}\times\frac{\sqrt{3}}{2}+\frac{3}{5}\times\frac{1}{2}=\frac{3+4\sqrt{3}}{10}\).

解法二：由 \(\tan \frac{\alpha}{2}=\frac{1}{2}, 0 < \frac{1}{2} < 1\)，知 \(\alpha\) 为第一象限角. 又
\[
\tan \alpha=\frac{2 \tan \frac{\alpha}{2}}{1-\tan^2 \frac{\alpha}{2}}=\frac{2 \times \frac{1}{2}}{1-\frac{1}{4}}=\frac{4}{3}
\]
\[
\begin{array}{l}
\text{所以} \cos \alpha=\frac{3}{5}, \sin \alpha=\frac{4}{5} \\
\sin \left(\alpha+\frac{\pi}{6}\right)=\sin \alpha \cos \frac{\pi}{6}+\cos \alpha \sin \frac{\pi}{6} \\
=\frac{4}{5} \times \frac{\sqrt{3}}{2}+\frac{3}{5} \times \frac{1}{2} \\
=\frac{3+4\sqrt{3}}{10}.
\end{array}
\]
\end{solution}

\begin{problem}
公园要建造一个圆形的喷水池. 在水池中央垂直于水面安装一个花形柱子 \(OA\). \(O\) 恰在水面中心，\(OA=1.25\) 米. 安置在柱子顶端 \(A\) 处的喷头向外喷水，水流在各个方向上沿形状相同的抛物线路径落下，且在过 \(OA\) 的任一平面上抛物线路径如图所示. 为使水流形状较为漂亮，设计成水流在到 \(OA\) 距离为 \(1\) 米处达到距水面最大高度 \(2.25\) 米. 如果不计其它因素，那么水池的半径至少要多少米，才能使喷出的水流不致落到池外

\end{problem}

\begin{solution}
如图建立直角坐标系，则水流所呈现的抛物线方程为
\[
y=a(x-1)^2+2.25
\]
由题意，点 \(A\) 的坐标为 \((0,1.25)\)，把 \(x=0\) 和 \(y=1.25\) 代入上述方程，得 \(a=-1\)，于是抛物线方程为
\[
y=-(x-1)^2+2.25
\]
令 \(y=0\)，得 \(-(x-1)^2+2.25=0\)，解得 \(x_1=2.5\)，\(x_2=-0.5\)（不合，舍去）. 答：水池半径至少要 \(2.5\) 米，才能使水流不落到池外.
\end{solution}

\begin{problem}
如图，在三棱柱 \(ABC-A'B'C'\) 中，四边形 \(A'ABB'\) 是菱形，四边形 \(BCC'B'\) 是矩形，\(C'B' \perp AB\). （1）求证：平面 \(CA'B \perp\) 平面 \(A'AB\)；（2）若 \(C'B'=3\)，\(AB=4\)，\(\angle ABB'=60^\circ\)，求 \(AC'\) 与平面 \(BCC'\) 所成角的大小（用反三角函数表示）.

\end{problem}

\begin{solution}
（1）证明：因为在三棱柱 \(ABC-A'B'C'\) 中，\(C'B' \parallel CB\).
\[
\text{所以} CB \perp AB
\]
又因为 \(CB \perp BB'\)，\(AB \cap BB'=B\)，所以 \(CB \perp\) 平面 \(A'AB\). 因为 \(CB \subset\) 平面 \(CA'B\)，因此平面 \(CA'B \perp\) 平面 \(A'AB\).

（2）由 \(C'B' \perp\) 平面 \(A'AB\)，得平面 \(A'AB \perp\) 平面 \(BCC'\). 过 \(A\) 作 \(AH \perp\) 平面 \(BCC'\)，\(H\) 为垂足，则 \(H\) 在 \(BB'\) 上. 连接 \(C'H\)，则 \(\angle AC'H\) 为 \(AC'\) 与平面 \(BCC'\) 所成的角. 连接 \(AB'\)，由四边形 \(A'ABB'\) 是菱形，\(\angle ABB'=60^\circ\)，可知 \(\triangle ABB'\) 为等边三角形，而 \(H\) 为 \(BB'\) 中点. 又 \(AB'=4\)，\(AH=2\sqrt{3}\)，于是在 Rt\(\triangle C'B'A\) 中，\(AC'=\sqrt{4^2+3^2}=5\)，而在 Rt\(\triangle AHC'\) 中，\(\sin \angle AC'H=\frac{2\sqrt{3}}{5}\)，
\[
\text{所以} \angle AC'H=\arcsin \frac{2\sqrt{3}}{5}
\]
因此，直线 \(AC'\) 与平面 \(BCC'\) 所成的角是 \(\arcsin \frac{2\sqrt{3}}{5}\).
\end{solution}

\section{选作题}

\begin{problem}
（\(A_2\)）设虚数 \(z_1\)，\(z_2\) 满足 \(z_1^2=z_2\). （1）若 \(z_1\)、\(z_2\) 又是一个实系数一元二次方程的两个根，求 \(z_1\)、\(z_2\)；（2）若 \(z_1=1+m\i\)（\(m>0\)，\(\i\) 为虚数单位），\(\omega=z_2-2\)，\(\omega\) 的辐角主值为 \(\theta\)，求 \(\theta\) 的取值范围.

\end{problem}

\begin{solution}
（1）解法一：设 \(z_1=a+bi\)，\(z_2=a-bi\)（\(a, b \in \R\) 且 \(b \neq 0\)）. 由 \(z_1^2=z_2\)，得 \(a^2-b^2+2abi=a-bi\)，于是
\[
\left\{
\begin{aligned}
a^2-b^2&=a,\\
2ab&=-b.
\end{aligned}
\right.
\]
解得
\[
\left\{
\begin{array}{l}
a=-\frac{1}{2},\\
b=\frac{\sqrt{3}}{2},
\end{array}
\right.
\]
\[
\text{所以} z_1=-\frac{1}{2}+\frac{\sqrt{3}}{2}\i,\ z_2=-\frac{1}{2}-\frac{\sqrt{3}}{2}\i
\]
或 \(z_1=-\frac{1}{2}-\frac{\sqrt{3}}{2}\i\)，\(z_2=-\frac{1}{2}+\frac{\sqrt{3}}{2}\i\).

解法二：设实系数一元二次方程为 \(z^2+pz+q=0\)（\(p^2-4q<0\)）. 则
\[
z_1=\frac{-p\pm\sqrt{4q-p^2}\i}{2}, \qquad z_2=\frac{-p\mp\sqrt{4q-p^2}\i}{2}
\]
由 \(z_1^2=z_2\)，得
\[
\frac{p^2-(4q-p^2)}{4} \pm \frac{-2p\sqrt{4q-p^2}\i}{4}=-\frac{p}{2} \mp \frac{\sqrt{4q-p^2}\i}{2}\i
\]
于是
\[
\left\{
\begin{array}{l}
\frac{p^2-2q}{2}=-\frac{p}{2},\\
\frac{-p\sqrt{4q-p^2}}{2}=-\frac{\sqrt{4q-p^2}}{2}.
\end{array}
\right.
\]
解得
\[
\left\{
\begin{array}{l}
q=1,
\end{array}
\right.
\]
\[
\text{所以} z_1=-\frac{1}{2}+\frac{\sqrt{3}}{2}\i,\ z_2=-\frac{1}{2}-\frac{\sqrt{3}}{2}\i
\]
或 \(z_1=-\frac{1}{2}-\frac{\sqrt{3}}{2}\i\)，\(z_2=-\frac{1}{2}+\frac{\sqrt{3}}{2}\i\).

（2）由 \(z_1=1+mi\)（\(m>0\)），\(z_1^2=z_2\)，得 \(z_2=(1-m^2)+2mi\)，
\[
\text{所以} \omega=-(1+m^2)+2m\i
\]
\[
\tan \theta=-\frac{2m}{1+m^2}=-\frac{2}{m+\frac{1}{m}}
\]
由 \(m>0\)，知 \(m+\frac{1}{m} \geqslant 2\). 于是 \(-1 \leqslant \tan \theta < 0\). 又由 \(-(m^2+1)<0\)，\(2m>0\)，得 \(\frac{3\pi}{4} \leqslant \theta < \pi\). 因此，所求 \(\theta\) 的取值范围为 \(\left[\frac{3\pi}{4}, \pi\right)\).
\end{solution}

\begin{problem}
（\(B_2\)）如图，抛物线 \(C_1:y=-x^2\) 与抛物线 \(C_2:y=x^2-2ax\)（\(a>0\)）交于 \(O\)、\(A\) 两点. （1）把 \(C_1\) 与 \(C_2\) 所围成的图形（阴影部分）绕 \(x\) 轴旋转一周，求所得几何体的体积 \(V\)；（2）若过原点的直线 \(l\) 与抛物线 \(C_2\) 所围成的图形面积为 \(\frac{9}{2}a^3\)，求直线 \(l\) 的方程.

\end{problem}

\begin{solution}
（1）\(C_1\) 和 \(C_2\) 交于两点 \(O(0,0)\)、\(A(a,-a^2)\)，
\[
\begin{array}{l}
V=\pi \int_0^a (x^2-2ax)^2 \mathrm{d}x-\pi \int_0^a (-x^2)^2 \mathrm{d}x \\
=4a\pi \int_0^a (-x^3+ax^2) \mathrm{d}x \\
=4a\pi \left(-\frac{1}{4}x^4+\frac{1}{3}ax^3\right)\Bigg|_0^a=\frac{\pi}{3}a^5.
\end{array}
\]

（2）设过原点的直线方程为 \(y=kx\)，解方程组
\[
\left\{
\begin{array}{l}
y=kx,\\
y=x^2-2ax,
\end{array}
\right.
\]
得 \(x_1=0\)，\(x_2=k+2a\). 当 \(k+2a \geqslant 0\) 时，
\[
\begin{array}{l}
S=\int_0^{k+2a} (kx-x^2+2ax) \mathrm{d}x \\
=\int_0^{k+2a} \left[(k+2a)x-x^2\right] \mathrm{d}x \\
=\left(\frac{k+2a}{2}x^2-\frac{1}{3}x^3\right)\Bigg|_0^{k+2a} \\
=\frac{(k+2a)^3}{6}.
\end{array}
\]
于是 \((k+2a)^3=27a^3\)，解得 \(k=a\). 所以，直线 \(l\) 的方程为 \(y=ax\).

当 \(k+2a<0\) 时，
\[
S=\int_{k+2a}^0 \left[(k+2a)x-x^2\right] \mathrm{d}x=-\frac{(k+2a)^3}{6}
\]
于是 \(-(k+2a)^3=27a^3\)，解得 \(k=-5a\). 所以，直线 \(l\) 的方程为 \(y=-5ax\).
\end{solution}

\section{解答题}

\begin{problem}
如图，抛物线方程为 \(y^2=p(x+1)\)（\(p>0\)），直线 \(x+y=m\) 与 \(x\) 轴的交点在抛物线的准线的右边.

\end{problem}

\begin{solution}
（1）证明：抛物线 \(y^2=p(x+1)\) 的准线方程是 \(x=-1-\frac{p}{4}\)，直线 \(x+y=m\) 与 \(x\) 轴的交点为 \((m,0)\)，则 \(m>-1-\frac{p}{4}\)，即 \(4m+p+4>0\).

由
\[
\left\{
\begin{aligned}
y^2&=p(x+1),\\
x+y&=m,
\end{aligned}
\right.
\]
得 \(x^2-(2m+p)x+(m^2-p)=0\)，而判别式 \(\Delta=(2m+p)^2-4(m^2-p)=p(4m+p+4)\). 又 \(p>0\) 及 \(4m+p+4>0\)，可见 \(\Delta>0\). 因此，直线与抛物线总有两个交点.

（2）设 \(Q\)、\(R\) 两点的坐标分别为 \((x_1,y_1)\)、\((x_2,y_2)\)，由（1）知，\(x_1\)，\(x_2\) 是方程 \(x^2-(2m+p)x+m^2-p=0\) 的两根. 所以 \(x_1+x_2=2m+p\)，\(x_1x_2=m^2-p\). 由 \(OQ \perp OR\)，得 \(KOQ \cdot KO_R=-1\)，即有 \(x_1x_2+y_1y_2=0\). 又 \(Q\)，\(R\) 为直线 \(x+y=m\) 上的点，因而
\[
y_1=-x_1+m, \qquad y_2=-x_2+m
\]
于是 \(x_1x_2+y_1y_2=2x_1x_2-m(x_1+x_2)+m^2=0\)，
\[
\text{所以} p=f(m)=\frac{m^2}{m+2}
\]
由
\[
\left\{
\begin{aligned}
p&>0,\\
4m+4+p&>0,
\end{aligned}
\right.
\]
得 \(m>-2\)，\(m \neq 0\).

（3）解法一：由于原点 \(O\) 到直线 \(x+y=m\) 的距离不大于 \(\frac{\sqrt{2}}{2}\)，于是
\[
\frac{|0+0-m|}{\sqrt{2}} \leqslant \frac{\sqrt{2}}{2}
\]
\[
\text{所以} |m| \leqslant 1
\]
由（2）知 \(m>-2\) 且 \(m \neq 0\)，故
\[
m \in [-1,0) \cup (0,1]
\]
由（2），
\[
f(m)=\frac{m^2}{m+2}=(m+2)+\frac{4}{m+2}-4
\]
当 \(m \in [-1,0)\) 时，任取 \(m_1\)，\(m_2\)，\(0>m_1>m_2 \geqslant -1\)，则
\[
f(m_1)-f(m_2)=(m_1-m_2)+\left(\frac{4}{m_1+2}-\frac{4}{m_2+2}\right)
\]
\[
=(m_1-m_2)\left[1-\frac{4}{(m_1+2)(m_2+2)}\right]
\]
由 \(0>m_1>m_2 \geqslant -1\)，知 \(0<(m_1+2)(m_2+2)<4\)，\(1-\frac{4}{(m_1+2)(m_2+2)}<0\). 又由 \(m_1-m_2>0\)，知 \(f(m_1)<f(m_2)\)，因而 \(f(m)\) 为减函数. 可见，当 \(m \in [-1,0)\) 时，\(p \in (0,1]\). 同样可证，当 \(m \in (0,1]\) 时，\(f(m)\) 为增函数，从而 \(p \in \left(0,\frac{1}{3}\right]\).

解法二：由解法一知，\(m \in [-1,0) \cup (0,1]\). 由（2）知，
\[
p=f(m)=\frac{m^2}{m+2}=\frac{1}{\frac{1}{m}+\frac{2}{m^2}}
\]
设 \(t=\frac{1}{m}\)，\(g(t)=t+2t^2\)，则 \(t \in (-\infty,-1] \cup [1,+\infty)\)，又
\[
g(t)=2t^2+t=2\left(t+\frac{1}{4}\right)^2-\frac{1}{8}
\]
因此当 \(t \in (-\infty,-1]\) 时，\(g(t)\) 为减函数，\(g(t) \in [1,+\infty)\)；当 \(t \in (1,+\infty)\) 时，\(g(t)\) 为增函数，\(g(t) \in [3,+\infty)\). 可见，当 \(m \in [-1,0)\) 时，\(t \in (-\infty,-1]\)，\(p=\frac{1}{g(t)} \in (0,1]\)；当 \(m \in (0,1]\) 时，\(t \in [1,+\infty)\)，\(p \in \left(0,\frac{1}{3}\right]\).
\end{solution}

\begin{problem}
设数列 \(\{a_n\}\) 的首项 \(a_1=1\)，前 \(n\) 项和 \(S_n\) 满足关系式：\(3tS_n-(2t+3)S_{n-1}=3t\)（\(t>0\)，\(n=2,3,4,\dots\)）. （1）求证：数列 \(\{a_n\}\) 是等比数列；（2）设数列 \(\{a_n\}\) 的公比为 \(f(t)\)，作数列 \(\{b_n\}\)，使 \(b_1=1\)，\(b_n=f\left(\frac{1}{b_{n-1}}\right)\)（\(n=2,3,4,\cdots\)），求 \(\displaystyle\lim_{n\to\infty}\frac{\lg a_n}{b_n}\)；（3）求和：\(b_1b_2-b_2b_3+b_3b_4-\cdots+(-1)^{n-1}b_nb_{n+1}\).

\end{problem}

\begin{solution}
（1）由 \(S_1=a_1=1\)，\(S_2=a_1+a_2=1+a_2\)，得
\[
3t(1+a_2)-(2t+3)=3t
\]
可见 \(a_2=\frac{2t+3}{3t}\)，\(\frac{a_2}{a_1}=\frac{2t+3}{3t}\). 又 \(3tS_n-(2t+3)S_{n-1}=3t\)，
\[
3tS_{n-1}-(2t+3)S_{n-2}=3t
\]
两式相减，得 \(3ta_n-(2t+3)a_{n-1}=0\). 于是 \(\frac{a_n}{a_{n-1}}=\frac{2t+3}{3t}\)（\(n=3,4,\dots\)）. 因此，\(\{a_n\}\) 是一个首项为 \(1\)，公比为 \(\frac{2t+3}{3t}\) 的等比数列.

（2）由 \(f(t)=\frac{2t+3}{3t}=\frac{2}{3}+\frac{1}{t}\)，\(b_n=f\left(\frac{1}{b_{n-1}}\right)=\frac{2}{3}+b_{n-1}\)，可见，\(\{b_n\}\) 是一个首项为 \(1\)，公差为 \(\frac{2}{3}\) 的等差数列. 于是 \(b_n=1+\frac{2}{3}(n-1)=\frac{2n+1}{3}\).
\[
\begin{array}{l}
\text{所以} \lim_{n \rightarrow \infty} \frac{\lg a_n}{b_n}=\lim_{n \rightarrow \infty} \frac{3(n-1)\lg \left(\frac{2t+3}{3t}\right)}{2n+1} \\
=\frac{3}{2} \lg \left(\frac{2t+3}{3t}\right).
\end{array}
\]

（3）由 \(b_n=\frac{2n+1}{3}\)，可知 \(\{b_{2n-1}\}\) 和 \(\{b_{2n}\}\) 是首项分别为 \(1\) 和 \(\frac{5}{3}\)，公差均为 \(\frac{4}{3}\) 的等差数列，于是 \(b_{2n}=\frac{4n+1}{3}\)，\(b_{2n+1}=\frac{4n+3}{3}\). 当 \(n=2m\)（\(m=1,2,\dots\)）时，
\[
\begin{array}{l}
b_1b_2-b_2b_3+b_3b_4-b_4b_5+\cdots+b_{2m-1}b_{2m}-b_{2m}b_{2m+1} \\
=b_2(b_1-b_3)+b_4(b_3-b_5)+\cdots+b_{2m}(b_{2m-1}-b_{2m+1}) \\
=-\frac{4}{3}(b_2+b_4+\cdots+b_{2m}) \\
=-\frac{4}{3} \cdot \frac{1}{2}m \cdot \left(\frac{5}{3}+\frac{4m+1}{3}\right) \\
=-\frac{4}{9}m(2m+3)=-\frac{1}{9}(2n^2+6n).
\end{array}
\]
当 \(n=2m-1\)（\(m=1,2,\dots\)）时，
\[
\begin{array}{l}
b_1b_2-b_2b_3+b_3b_4-b_4b_5+\cdots-b_{2m-2}b_{2m-1}+b_{2m-1}b_{2m} \\
=-\frac{4}{9}m(2m+3)+b_{2m}b_{2m+1} \\
=-\frac{4}{9}m(2m+3)+\frac{(4m+1)(4m+3)}{9} \\
=\frac{1}{9}(8m^2+4m+3)=\frac{1}{9}(2n^2+6n+7).
\end{array}
\]
\[
\begin{array}{l}
\text{所以} b_1b_2-b_2b_3+b_3b_4-\dots+(-1)^{n-1}b_nb_{n+1} \\
=\left\{
\begin{array}{ll}
-\frac{1}{9}(2n^2+6n), & n \text{为偶数}; \\
\frac{1}{9}(2n^2+6n+7), & n \text{为奇数}.
\end{array}
\right.
\end{array}
\]
\end{solution}
